% !TeX root = ../../book.tex
本章的目标是发展一种有条不紊的方法，将一个命题分解成更小的组成部分，并了解这些组成部分如何组合在一起---这被称为命题的\textit{逻辑结构}。一个命题的逻辑结构非常有用：它告诉我们需要做什么来证明它，需要写什么来表达我们的证明，以及在命题被证明后如何探索命题的后续结果。

\begin{center} \vspace{-15pt}
\fitwidth{%
\begin{tikzcd}[row sep={30pt}, column sep={20pt}, ampersand replacement=\&]
\&
\text{\begin{minipage}{90pt}\centering 命题的逻辑结构\end{minipage}}
\arrow[dl]
\arrow[d]
\arrow[dr]
\&
\\
\text{\begin{minipage}{120pt}\centering 命题的证明策略\end{minipage}}
\&
\text{\begin{minipage}{100pt}\centering 证明的结构和措辞 \end{minipage}}
\&
\text{\begin{minipage}{100pt}\centering 命题的后续结果\end{minipage}}
\end{tikzcd}%
}
\end{center}

\Cref{secPropositionalLogic,secVariablesQuantifiers} 致力于发展用于命题推理的\textit{符号逻辑}系统。我们可以使用一串变量和符号来表示一个命题，这个表达式将指导我们如何证明该命题并探索其结果。在 \Cref{secLogicalEquivalence} 中，我们将发展起以代数方式处理这些逻辑表达式的技术---这反过来会产生新的证明技术（有些有奇特的名称，如`对位证明'，有些则没有）。

探索命题的逻辑结构如何影响其证明的结构和措辞是 \Cref{secVocabulary} 的内容。
